\begin{abstract}
    Co-training, which combines limited in-domain real-world data with abundant surrogate data such as simulation or cross-embodiment robot data, is widely used for training generative robot policies. Despite its empirical success, the mechanisms that determine when and why co-training is effective remain poorly understood.
    We investigate the mechanism of sim-and-real co-training through theoretical analysis and empirical study, and identify two intrinsic effects governing performance. The first, \textbf{``structured representation alignment"}, reflects a balance between cross-domain representation alignment and domain discernibility, and plays a primary role in downstream performance. The second, the \textbf{``importance reweighting effect"}, arises from domain-dependent modulation of action weighting and operates at a secondary level.
    We validate these effects with controlled experiments on a toy model and extensive sim-and-sim and sim-and-real robot manipulation experiments. Our analysis offers a unified interpretation of recent co-training techniques and motivates a simple method that consistently improves upon prior approaches.
    More broadly, our aim is to examine the inner workings of co-training and to facilitate research in this direction.
\end{abstract}
\vspace{-1em}